% Generated draft: requires mathematical and source review.
\documentclass[11pt]{article}
\usepackage[margin=25mm]{geometry}
\usepackage{fontspec}
\setmainfont{Latin Modern Roman}
\usepackage{amsmath,amssymb,mathtools,unicode-math}
\setmathfont{Latin Modern Math}
\usepackage{enumitem,needspace,xurl,hyperref}
\setlength{\parindent}{0pt}
\setlength{\parskip}{5pt}
\setlength{\emergencystretch}{3em}
\newcommand{\N}{\mathbb{N}}
\newcommand{\Z}{\mathbb{Z}}
\newcommand{\Q}{\mathbb{Q}}
\newcommand{\R}{\mathbb{R}}
\newcommand{\C}{\mathbb{C}}
\newcommand{\F}{\mathbb{F}}
\newcommand{\A}{\mathbb{A}}
\newcommand{\PP}{\mathbb{P}}
\newcommand{\E}{\mathbb{E}}
\newcommand{\eps}{\varepsilon}
\newcommand{\dd}{\,\mathrm{d}}
\begin{document}
\section*{UnsolvedMath ID: 20002093}
\textbf{Source problem code:} catalogue entry
\subsection{Definitions and mathematical statement}
The aid's source-boundary correction says numbered Question 1.35 consists only of ``Eliashberg made some remarks on the above topics.'' It contains no unresolved mathematical question. The following paragraphs are Banyaga's unnumbered exposition on locally conformal symplectic geometry; Questions 1.36 and 1.37 are separate records. The registry's question/open metadata is retained without inventing an open conjecture here.

In the source's sign convention, an l.c.s. structure is a nondegenerate two-form \(\Omega\) with a closed one-form \(\eta\) satisfying
\[
 \dd\Omega=-\eta\wedge\Omega.
\]
For closed \(\eta\), define \(\dd_\eta\beta=\dd\beta+\eta\wedge\beta\), so \(\dd_\eta^2=0\), \(\dd_\eta\Omega=0\), and
\[
 H_\eta^k(X)=H^k(\Omega^\bullet(X),\dd_\eta)
\]
is the twisted/Lichnerowicz cohomology in this convention. Renaming the Lee form avoids the source's reuse of \(\omega\) for both a one-form and a symplectic form.

For a contact manifold \((N,\alpha)\), on \(N\times S^1\) let \(\vartheta\) be the pulled-back contact form and \(\eta\) a pulled-back nowhere-zero volume one-form on the circle. Then
\(\Omega=\dd\vartheta+\eta\wedge\vartheta=\dd_\eta\vartheta\)
is the source's l.c.s. example. Another is \(\Omega=f\sigma\), where \(\sigma\) is symplectic and \(f>0\), with \(\eta=-\dd\log f\); this is globally conformal symplectic. Conversely, an exact Lee form in this sign convention can be removed by a positive global conformal factor.

The aid develops a strict-contact mapping-torus extension. Let \(N\) be a compact contact manifold without boundary, \(\phi:N\to N\) a strict contactomorphism \(\phi^*\alpha=\alpha\), and \(a\in\R\setminus\{0\}\). Define
\[
 X_\phi=(N\times[0,1])/((x,1)\sim(\phi(x),0)),
 \qquad \eta=a\,\dd t.
\]
The form \(\vartheta\) induced by \(\alpha\) descends, and
\[
 \Omega=\dd_\eta\vartheta=\dd\vartheta+a\dd t\wedge\vartheta
\]
is nondegenerate and l.c.s.; its Lee period on the base circle is \(a\).

Over \(\R\), put \(A_j=e^a\phi^*-I\) on \(H^j(N;\R)\). The aid reports the Wang short exact sequence
\[
 0\longrightarrow\operatorname{coker}A_{k-1}
 \longrightarrow H_\eta^k(X_\phi)
 \longrightarrow\ker A_k\longrightarrow0.
\]
No canonical splitting is asserted. If
\[
 m_j(e^{-a})=\dim_\R\ker(\phi^*-e^{-a}I:H^j(N;\R)\to H^j(N;\R)),
\]
with cohomology outside the degree range zero, then
\[
 \dim H_\eta^k(X_\phi)=m_k(e^{-a})+m_{k-1}(e^{-a}).
\]
These are geometric eigenspace multiplicities, not algebraic multiplicities of Jordan blocks. All twisted groups vanish exactly when \(e^{-a}\) is absent from the cohomological monodromy spectrum in every degree. In the product case \(\phi=\operatorname{id}\) and \(a\ne0\), the criterion gives vanishing. This is the aid's developed contextual result, not a solution of a missing numbered question.
\subsection{Short English statement}
The numbered item is only a historical remark; the l.c.s. definitions belong to following exposition. The aid extends the contact-product example to strict-contact mapping tori and computes twisted cohomology from monodromy eigenspaces. Its sign and Lee-period conventions are explicit, and geometric multiplicity is kept distinct from algebraic multiplicity.
\subsection{Sources}
\begin{itemize}
\item Local UnsolvedMath registry, permanent ID 20002093, exact source code catalogue entry; canonical \texttt{aim-geometry-notes.json}, \texttt{notes[430]}, tag \texttt{question}; AIM Workshop Problem Lists, Geometry, \emph{Holomorphic curves in contact geometry}. Created and updated 14 August 2026.
\item Source URL \url{https://aimath.org/WWN/contactgeom2/contactgeom2.pdf}. The supplied verification and aid distinguish the Eliashberg remark numbered 1.35 from Banyaga's following exposition and the separately numbered 1.36--1.37. The official PDF/HTML layout is reported as checked in the supplied evidence; neither was accessed here.
\item Lichnerowicz cohomology is named in the source. No separate authoritative citation is supplied for the mapping-torus Wang calculation or its spectral criterion.
\item Attempt 1: \texttt{context\_only}, problem status at run \texttt{not\_a\_problem}, rigor label \texttt{proved}, novelty confidence \texttt{low}; classification \texttt{OPEN-TRIAGE}, suggested difficulty L2, registry status \texttt{open}. Candidate: strict-contact mapping-torus spectral criterion. These labels are retained without treating context as an independent conjecture.
\item Ulam AI / UnsolvedMath Contributors; CC BY 4.0. Machine-generated research aid, not peer reviewed; independently verify before citation.
\end{itemize}
\end{document}
